\section{Introduction}
A fixed per-image channel-symbol budget limits the source description available for perceptually faithful reconstruction of image structure, appearance, and content. Its useful content depends on image-dependent source-code length and the protection carried in the same frame. A generative receiver can complete untransmitted information, while the received description constrains that completion to the source image. The challenge is to retain useful source information within the frame and use it to improve perceptual similarity to the source.

Next-scale representations organize an image into complete coarse token maps followed by residual refinements. After resizing, each coarse map covers the full latent grid, providing a global description before finer maps add detail. Visual autoregressive modeling (VAR) reported state-of-the-art results in its original class-conditional ImageNet generation benchmarks~\cite{var}, motivating its use as a completion prior. Digital coding and learned joint source--channel coding (JSCC) address limited channel resources~\cite{deepjscc,swinjscc}; adapting progressive descriptions is a classical problem~\cite{chande}. Recent token schemes include Ada-TokenCom's prefix coding and tail generation~\cite{adatokencom}, task-oriented TONIC~\cite{tonic}, and next-scale progressive ARPC~\cite{arpc}.

Whole-scale boundaries need not align with fixed frame capacity: a scale adds many tokens, and the terminated entropy-code length varies across images. Legal endpoints inside a scale can retain source tokens that would otherwise be excluded. VAR's coarse-conditioned factorization lets both terminals construct matching probabilities for these positions from complete coarse maps. The delivered boundary tokens must then refine both residual accumulation and the context for predicting finer scales.

NeST-Com combines within-scale endpoints, independently terminated arithmetic streams, actual-length fitting, and prefix-constrained reconstruction. Using VAR's known-token overwrite~\cite{var}, it merges received positions with predictions and updates both the residual latent and finer-scale context. With shared model components on common500, it reduces LPIPS and DreamSim at all six tested budget/SNR points, with PSNR gains of 0.32--0.79~dB over whole-scale transmission. At 1024 symbols and 10~dB, source-bit utilization rises from 77.70\% to 95.21\%, while mean LPIPS and DreamSim fall by 13.7\% and 18.3\%. Controlled experiments examine how retained boundary information affects reconstruction; complete-system comparisons reveal perceptual and pixel fidelity tradeoffs.
